\subsection{\texorpdfstring{Condition for \(\mathcal{L}_{\mathfrak{S}}(\mathbf{E};\mathbf{F})\) to be Hausdorff}{Condition for \textbackslash mathcal\{L\}\_\{\textbackslash mathfrak\{S\}\}(\textbackslash mathbf\{E\};\textbackslash mathbf\{F\}) to be Hausdorff}}

PROPOSITION 3. -- Let E and F be two locally convex spaces, F being assumed Hausdorff, and let \(\mathfrak{S}\) be a family of bounded subsets of E. If the union A of the sets of \(\mathfrak{S}\) is total in E, then the space \(\mathcal{L}_{\mathfrak{S}}(\mathrm{E};\mathrm{F})\) is Hausdorff.

Let \(u_0\) be a non-zero element of \(\mathcal{L}(\mathrm{E};\mathrm{F})\) ; since \(u_0\) is continuous and \(\mathrm{A}\) is total in \(\mathrm{E}\) , there exists an \(x_0\) in \(\mathrm{A}\) such that \(u_0(x_0) \neq 0\) . Since \(\mathrm{F}\) is Hausdorff, there exists a neighbourhood \(\mathrm{V}\) of 0 in \(\mathrm{F}\) such that \(u_0(x_0) \notin \mathrm{V}\) . Let \(\mathbf{M} \in \mathfrak{S}\) be such that \(x_0 \in \mathbf{M}\) . Then the set \(\mathrm{U}\) of all \(u \in \mathcal{L}(\mathrm{E};\mathrm{F})\) such that \(u(\mathbf{M}) \subset \mathbf{V}\) is a neighbourhood of 0 in \(\mathcal{L}(\mathrm{E};\mathrm{F})\) , and we have \(u_0 \notin \mathrm{U}\) , hence \(\mathcal{L}(\mathrm{E};\mathrm{F})\) is Hausdorff.

In particular, the following topologies on \(\mathcal{L}(\mathrm{E};\mathrm{F})\) are Hausdorff whenever F is Hausdorff: simple convergence, compact convergence, precompact or compact convex, and bounded convergence.

